\par\Needspace{6\baselineskip}\medskip
\begin{center}\hypertarget{native-243-article173}{}\pdfbookmark[2]{173.}{article-173}173.\end{center}
Besonders wichtig ist es, jede geozentrische Position des Himmelskörpers, auf der die Bahnbestimmung beruhen soll, möglichst aus mehreren Beobachtungen zu gewinnen, die so kombiniert werden, dass sich etwaige Fehler so weit wie möglich gegenseitig aufheben, statt sie aus einer einzelnen Beobachtung abzuleiten. Beobachtungen, die nur wenige Tage auseinanderliegen --- oder je nach Lage der Dinge auch 15 oder 20 Tage ---, gehen also nicht als ebenso viele verschiedene Positionen in die Rechnung ein. Vielmehr wird daraus eine einzige Position abgeleitet, die gewissermaßen das Mittel aller Beobachtungen bildet und daher eine weit höhere Präzision zulässt als die einzelnen Beobachtungen für sich genommen. Dieses Verfahren beruht auf den folgenden Grundsätzen.

Die aus den näherungsweise bekannten Bahnelementen berechneten geozentrischen Orte des Himmelskörpers dürften nur wenig von den wahren Orten abweichen. Die Differenzen zwischen ihnen dürften sich nur sehr langsam ändern, sodass sie im Laufe weniger Tage als nahezu konstant gelten können oder zumindest ihre Änderungen als proportional zu den Zeiten anzusehen sind. Wären also die Beobachtungen völlig fehlerfrei, so wären die Differenzen zwischen den zu den Zeiten $t$, $t'$, $t''$, $t'''$ usw. beobachteten Orten und den aus den Bahnelementen berechneten Orten, d.\,h. die Differenzen zwischen beobachteten und berechneten Längen und Breiten beziehungsweise Rektaszensionen und Deklinationen, entweder praktisch gleich oder sie würden zumindest gleichmäßig und sehr langsam zu- oder abnehmen. Seien z.\,B. die zu diesen Zeiten beobachteten Rektaszensionen $a$, $a'$, $a''$, $a'''$ usw. und die berechneten $a+\delta$, $a'+\delta'$, $a''+\delta''$, $a'''+\delta'''$ usw.; dann weichen die Differenzen $\delta$, $\delta'$, $\delta''$, $\delta'''$ usw. von den wahren, durch die Bahnelemente bedingten Abweichungen nur insoweit ab, als die Beobachtungen selbst fehlerhaft sind: Lassen sich also jene Abweichungen für alle diese Beobachtungen als konstant ansehen, so liefern die Größen $\delta$, $\delta'$, $\delta''$, $\delta'''$ usw. ebenso viele verschiedene Bestimmungen derselben Größe. Als ihren richtigen Wert wird man daher das arithmetische Mittel dieser Bestimmungen ansetzen, sofern kein Grund besteht, die eine \SourceLeaf{244}oder die andere vorzuziehen. Scheint es jedoch nicht angebracht, allen Beobachtungen denselben Grad an Präzision zuzuschreiben, so nehmen wir an, dass die Präzision der einzelnen Beobachtungen jeweils als proportional zu den Zahlen $e$, $e'$, $e''$, $e'''$ usw. zu veranschlagen ist, d.\,h. dass bei den Beobachtungen Fehler, die zu diesen Zahlen umgekehrt proportional sind, gleich leicht auftreten konnten; dann ist nach den weiter unten darzulegenden Grundsätzen der wahrscheinlichste Mittelwert nicht mehr das einfache arithmetische Mittel, sondern
\[
=\frac{ee\delta+e'e'\delta'+e''e''\delta''+e'''e'''\delta'''+\text{usw.}}
{ee+e'e'+e''e''+e'''e'''+\text{usw.}}\mathpunct{\cdot}
\]
Setzen wir nun diesen Mittelwert $=\Delta$, so können wir als wahre Rektaszensionen der Reihe nach $a+\delta-\Delta$, $a'+\delta'-\Delta$, $a''+\delta''-\Delta$, $a'''+\delta'''-\Delta$ usw. ansetzen, und es bleibt uns überlassen, welche davon wir in der Rechnung verwenden. Liegen jedoch die Beobachtungen zeitlich zu weit auseinander oder sind die Bahnelemente noch nicht hinreichend genau bekannt, sodass sich ihre Abweichungen nicht für alle Beobachtungen als konstant ansehen lassen, so ergibt sich offenbar nur der Unterschied, dass die so gefundene mittlere Abweichung auf eine bestimmte mittlere Zeit zu beziehen ist, statt sie für alle Beobachtungen als dieselbe anzunehmen. Die mittlere Zeit muss aus den einzelnen Zeitpunkten ebenso hergeleitet werden wie $\Delta$ aus den einzelnen Abweichungen. Allgemein ergibt sich also die Zeit
$\frac{eet+e'e't'+e''e''t''+e'''e'''t'''+\text{usw.}}{ee+e'e'+e''e''+e'''e'''+\text{usw.}}\mathpunct{\cdot}$
Strebt man die höchste Präzision an, so muss man für diese Zeit den geozentrischen Ort aus den Bahnelementen berechnen und anschließend den mittleren Fehler $\Delta$ abziehen, um eine möglichst genaue Position zu erhalten: In der Regel genügt es jedoch vollkommen, den mittleren Fehler auf die Beobachtung zu beziehen, die der mittleren Zeit am nächsten liegt. Was hier über die Rektaszensionen gesagt wurde, gilt ebenso für die Deklinationen oder auch, wenn man es vorzieht, für die Längen und Breiten: Es ist jedoch stets günstiger, die aus den Bahnelementen berechneten Rektaszensionen und Deklinationen unmittelbar mit den beobachteten zu vergleichen; damit vereinfacht sich die Rechnung, besonders wenn man die in Art.\ 53---60 dargelegten Verfahren anwendet. Außerdem hat dieses Vorgehen den Vorzug, dass sich auch unvollständige Beobachtungen verwenden lassen. Würde man darüber hinaus alles auf Längen und Breiten beziehen, so wäre zu befürchten, dass eine hinsichtlich der Rektaszension gut, hinsichtlich der Deklination schlecht ausgeführte Beobachtung (oder umgekehrt) in beiden Koordinaten beeinträchtigt und damit völlig unbrauchbar würde. --- Der Grad an Präzision, der dem so gefundenen Mittel nach den sogleich zu erläuternden Grundsätzen zukommt, ist $=\surd(ee+e'e'+e''e''+e'''e'''+\text{usw.})$, sodass vier oder neun gleich \SourceLeaf{245}genaue Beobachtungen erforderlich sind, wenn das Mittel die doppelte oder dreifache Präzision haben soll, und so weiter.
